\documentclass[12pt,a4paper]{article}
\usepackage{ra_preamble}
\usepackage{longtable}
\title{\textbf{Relational Actualism: Irreversible Events\\
as the Primitive Basis of Physical Reality}\\[0.5em]
\large\normalfont Relational Actualism --- Paper~V: The Complete Framework}
\author{Joshua F.~Sandeman\\\small Independent Researcher, Salem, Oregon\\\small \texttt{sansuikyo@gmail.com} $\cdot$ \texttt{relationalactualism.org}}
\date{April 2026}
\begin{document}
\maketitle

\begin{abstract}
Relational Actualism (RA) is a framework in which irreversible on-shell actualization events --- physical interactions producing a permanent increase in quantum relative entropy $\Delta S(\rho\|\sigma_0)>\Delta S^*=0.60069\,\mathrm{nats}$ --- are the primitive physical reality.  Spacetime is the growing directed acyclic graph (DAG) of these events; time is the process of graph growth; the arrow of time is encoded in the irreversibility of actualization itself.

From a single ontological commitment, without free parameters, RA derives: $\alpha_{\rm EM}^{-1}=137$ \statusLV; $\alpha_s(m_Z)=1/\sqrt{72}$ \statusCV; Koide $K=2/3$ \statusLV; complete Standard Model spectrum \statusLV; $\Lambda=0$ \statusDR; BMV null prediction \statusDR; Hubble tension $H_0^{\rm Eridanus}\approx75.9\,\mathrm{km/s/Mpc}$ \statusCV; WIMP prohibition \statusDR; Boltzmann Brain prohibition \statusDR; minimum viable universe seed $N_{\min}=2$ \statusDR; universal decoherence timescale $\tau_d=20\,\mathrm{fs}$ at $300\,\mathrm{K}$ \statusDR; Kinematic Coherence Bound for quantum computing \statusAR; Causal Firewall origin-of-life threshold (same $N_{\min}=2$ as Big Bang) \statusDR.

Approximately 110 results are Lean~4 verified with zero sorry tags in the core files.  All 76 KB claims are documented with explicit epistemic status.  This paper is the hub of the series: Papers~I--IV are self-contained derivations; Paper~V provides the complete epistemic map, the coherence structure analysis, all predictions, all open problems, and the philosophical synthesis.
\end{abstract}

\tableofcontents\bigskip

\section{The Founding Commitment}\label{sec:commitment}

\subsection{The Ontological Claim}

\begin{mdframed}[linecolor=cv_blue!60,backgroundcolor=cv_blue!5,linewidth=2pt]
\textbf{The Relational Actualism Commitment:}\\[0.3em]
An interaction constitutes a \emph{physical event} if and only if it produces an irreversible increase in quantum relative entropy:
\begin{equation}
  \Delta S(\rho\|\sigma_0) > \Delta S^* = 0.60069\,\mathrm{nats}
\end{equation}
Such events write permanent directed edges into the growing causal DAG.  The DAG is reality.  Everything else is derived from it.
\end{mdframed}

This is not a hypothesis about the nature of spacetime alongside other hypotheses.  It is a proposal that irreversibility is the most fundamental fact about physical reality.  The universe is not a stage on which events play out; the universe \emph{is} the events.

\subsection{Time as the Growing Graph}

In standard physics, time is a parameter.  Events are \emph{located in} time.  The laws of physics are equations that describe how quantities evolve \emph{through} time.  Time is the background.

In RA, this is inverted.  Time \emph{is} the growing of the causal graph.  The past is the part of the graph that has been written.  The future is what has not yet been written.  The present is the process of the next edge being added.

Proper time along a worldline --- the time measured by a clock --- is the integer count of actualization events along that path.  Relativistic time dilation is the statement that different worldlines accumulate events at different rates.  The Lorentz factor $\gamma=dt/d\tau$ is the ratio of coordinate events to proper events.

The arrow of time is not explained by RA --- it is what RA \emph{is}.  Actualization is by definition irreversible.  The causal graph can only grow, never shrink.  There is no equation of motion in RA that runs both forward and backward; the BDG criterion $S(v,C)>0$ is not symmetric under time reversal.

\subsection{The BDG Action}

The dynamics are governed by the Benincasa--Dowker--Glaser action \cite{Dowker2013}:
\begin{equation}
  S(v,C) = 1 - N_1 + 9N_2 - 16N_3 + 8N_4
  \label{eq:BDG}
\end{equation}
The integers $(1,-1,9,-16,8)$ are uniquely determined by the requirement that this operator converges to the d'Alembertian (Paper~I, Section~2).  They are not free parameters.

\section{The Epistemic Map}\label{sec:epistemic}

\subsection{Why Explicit Status Matters}

The epistemic map is the most important methodological innovation in this series.  Theoretical physics has a history of overconfident claims: frameworks that were presented as derived when key steps were argued, or as established when they were conjectured.  RA attempts to reverse this tendency by documenting the epistemic status of every claim.

\begin{table}[ht]
\centering
\caption{Epistemic status categories.}
\label{tab:status}
\begin{tabular}{clp{9.5cm}}
\toprule
Badge & Name & Definition \\
\midrule
\statusLV & Lean-Verified & Machine-checked in Lean~4.  Zero sorry tags in core files.  Certainty as high as formal mathematics allows.\\
\statusCV & Computation-Verified & Certified numerical computation with explicit error bounds.  Reproducible by running the public scripts.\\
\statusDR & Derived & All steps explicit in the RA papers.  Argument is complete; the result follows from the premises.  Not machine-checked.\\
\statusAR & Argued & Key steps identified and physically motivated.  Some steps not fully explicit.  Upgradeable with further work.\\
\statusOP & Open & Precisely scoped gap with a stated path to resolution.\\
\bottomrule
\end{tabular}
\end{table}

Status does not measure importance.  Some of the most important results are \statusAR\ or \statusDR; some \statusLV\ results are supporting lemmas.  Status measures epistemic certainty.

\subsection{Complete Lean~4 Verification Table}

\begin{longtable}{llp{9cm}}
\caption{All Lean~4 verified results.  Files: RA\_D1\_Proofs.lean (main), RA\_AmpLocality.lean (amplitude locality), RA\_Spin2\_Macro.lean (spin-2), RA\_AQFT\_Proofs\_v10.lean (AQFT).  One intentional sorry in a library adapter unrelated to any RA claim.  Total: $\approx110$ verified results.}\label{tab:lean}\\
\toprule
ID & Theorem & Statement\\
\midrule
\endfirsthead
\multicolumn{3}{l}{\textit{(Table \ref{tab:lean} continued)}}\\
\toprule
ID & Theorem & Statement\\
\midrule
\endhead
\midrule\multicolumn{3}{r}{\textit{continued on next page}}\\
\endfoot
\bottomrule
\endlastfoot
L01 & LLC & $\Sigma_{\rm out}(v)=\Sigma_{\rm in}(v)$ at every DAG vertex \\
L02 & Graph Cut & LLC holds on both sides of any causal severance \\
L03 & Markov Blanket & Interior subgraph shielded from exterior \\
L04 & Frame Independence & $S(U\rho U^\dagger\|\sigma_0)=S(\rho\|\sigma_0)$ for all Poincar\'e $U$ \\
L05 & Rindler Stationarity & $\frac{d}{dt}S(\alpha_t(\rho_\beta)\|\sigma_0)=0$ under modular flow \\
L06 & Rindler Thermal & $\rho_\beta$ is valid density matrix (Hermitian, positive, $\mathrm{tr}=1$) \\
L07 & Causal Inv.\ (cond.) & $\mu_\sigma(S)=\mu_{\sigma'}(S)$ given amplitude locality \\
L08 & $\alpha_{\rm EM}$ & $144-7=137$ (\texttt{norm\_num}) \\
L09 & Koide & $K=(m_e+m_\mu+m_\tau)/(\sqrt{m_e}+\sqrt{m_\mu}+\sqrt{m_\tau})^2=2/3$ \\
L10 & Confinement depths & $L_{\rm gluon}=3$, $L_{\rm quark}=4$ \\
L11 & BDG Closure & 5 topology types, 124 extensions verified \\
L12 & Qubit Fragility & Min BDG score for electrons/photons $=1$ \\
O01 & Amplitude Locality & $a(v|C)$ depends only on $C\cap\mathrm{past}(v)$ \\
O02 & Causal Invariance & $\mu_\sigma(S)$ invariant for any permutation \\
O10s & Spin-2 DOF & $5-1-1-1=2$ physical degrees of freedom for BDG metric excitation \\
\end{longtable}

\subsection{Computation-Verified Results}

\begin{table}[ht]
\centering
\caption{All computation-verified results.  Scripts in the GitHub repository are public and reproducible.}
\label{tab:cv}
\begin{tabular}{llll}
\toprule
ID & Result & Value & Reference \\
\midrule
C01 & $\alpha_s(m_Z)=1/\sqrt{72}$ & $0.11785$ (PDG: $0.1180$) & d3\_alpha\_s\_proof.py \\
C02 & $P_{\rm acc}(1)$ & $0.548436$ & BDG Monte Carlo \\
C03 & $\Delta S^*$ & $0.60069$ nats (irrational) & exact enumeration \\
C04 & Spin-bath & $t^*=0.274/g$ & spin-bath model \\
C05 & Hubble tension & $H_0^{\rm Eridanus}=75.9\,\mathrm{km/s/Mpc}$ & void geometry \\
C06 & Proton mass & $m_p=\sqrt{8}\Lambda_{\rm QCD}$ (0.08\%) & BDG Regge \\
C07 & Roper resonance & $290\,\mathrm{MeV}$ gap = Kind~2 virtual & ChPT loop \\
C08 & Baryon fraction & $f_0=5.42$ (Planck: $5.416$) & BDG path weights \\
D4U02a & D4U02 analytic & $P(S=1)>P(S=0)$ by 3100:1 margin & modular arithmetic \\
D4U02u & $d=4$ unique & Modular barrier fails for $d=2,3$ & counting argument \\
\bottomrule
\end{tabular}
\end{table}

\section{The Coherence Hub}\label{sec:hub}

\subsection{Three Hub Results from April 2026}

In April 2026, three previously open problems were closed in a single session.  An audit of the new results revealed that each was a hub: each closed result enabled the closure or substantial progress of multiple other open problems.  This section documents the hub structure.

\paragraph{Hub 1: D4U02-analytic.}  The analytic proof that $P(S=1)>P(S=0)$ via conditioning and modular arithmetic:
\begin{itemize}
  \item Proved D4U02 analytically (OP$\to$DR) after six years.
  \item Established dimensional uniqueness ($d=4$ from modular barrier, new DR result).
  \item Opened new perturbative path for O03 (Bianchi on curved background).
  \item Motivated the $l_{\rm RA}=\ell_P$ identification from the dominant $k=2$ term (O06).
  \item Provided a concrete question for Dowker on Lorentz emergence (O11).
\end{itemize}

\paragraph{Hub 2: Universal timescale $\tau_{\rm act}=\Delta S^*/\Gamma_{\rm on\text{-}shell}$.}  The formula and its three-scale application:
\begin{itemize}
  \item Closed O07 (BMV timescale) as a derived result.
  \item Closed O05 (Unruh detector) via L05+timescale.
  \item Closed O08 (biological $\tau_d$) via timescale at $T=300\,\mathrm{K}$.
  \item Upgraded P01 (BMV null) from AR to DR.
  \item Connected cosmological and biological timescales through one formula.
\end{itemize}

\paragraph{Hub 3: COSMO-SEED bootstrap.}  The $N_{\min}=2$ condition:
\begin{itemize}
  \item Established exact BDG arithmetic for Big Bang seeds.
  \item Identified $N_{\min}=2$ = Causal Firewall origin-of-life threshold (same condition, two scales).
  \item Provided the exact Cosmological Natural Selection mechanism.
  \item Derived the BH mass hierarchy for viable universe production.
\end{itemize}

\subsection{The Universal Timescale as the Central Hub}

The formula $\tau_{\rm act}=\Delta S^*/(\Gamma_{\rm on\text{-}shell}\cdot\Delta S_{\rm per\,boson})$ in five physical contexts:
\begin{align}
  T=1\,\mathrm{mK}:& \quad\tau_{\rm act}\approx1.9\times10^4\,\mathrm{s} \gg \tau_{\rm exp} \quad\text{(BMV null)} \\
  T=300\,\mathrm{K}:& \quad\tau_d\approx20\,\mathrm{fs}=\tau_{\rm vib} \quad\text{(classical chemistry)} \\
  \text{Rindler}:& \quad\tau_{\rm act}\to\infty \quad\text{(KMS: $\Delta S_{\rm per\,photon}=0$)} \\
  N=2\text{ seed}:& \quad\tau_d=0 \quad\text{(instantaneous bootstrap)} \\
  \text{Spin-bath}:& \quad t^*=0.274/g \quad\text{(parameter-free)}
\end{align}
Fourteen orders of magnitude.  No free parameters.  One formula.

\subsection{The $N_{\min}=2$ Identity in Full}

The minimum viable universe seed (Paper~III) and the Causal Firewall origin-of-life threshold (Paper~IV) are the same condition: $N\geq2$ causally connected vertices with a directed edge.  BDG viability $V=1-N_1+9N_2=1-1+9=+9$ in both cases.

These are not coincidences.  They are the same geometric fact --- the minimum non-Markovian causal structure in the Poisson-CSG model --- applied at two different scales.  The Big Bang and the first self-replicating molecule share the same fundamental causal requirement.

\section{General Relativity and Cosmology}\label{sec:gr}

\subsection{From LLC to Einstein Equations}

The route to $G_{\mu\nu}=8\pi G\Pact[T_{\mu\nu}]$ with $\Lambda=0$ uses four Lean-verified inputs:
\begin{enumerate}
  \item L01 (LLC) gives conservation at every vertex.
  \item In the flat/weak-field limit (D04, \statusDR), LLC $\Rightarrow \nabla_\mu G^{\mu\nu}\equiv0$.
  \item L11 ($d=4$ from SM closure) activates Lovelock's uniqueness theorem.
  \item D03 (vacuum suppression) sets $\Lambda=0$.
\end{enumerate}
Combined with O10-spin2 (\statusLV: 2 physical DOF for the BDG metric excitation) and amplitude locality (O01, \statusLV), the Weinberg route \cite{Weinberg1964} gives the unique coupling: $G_{\mu\nu}=8\pi G\Pact[T_{\mu\nu}]$.

The cosmological constant $\Lambda=0$ because $\Pact[\rho_{\rm vac}]=0$: the vacuum consists of off-shell virtual processes with $\Delta S=0$, all below the $\Delta S^*$ threshold.

\subsection{Black Hole Thermodynamics}

$S_{\rm BH}=A/(4\ell_P^2)$ follows from D10 (\statusAR, Bridge Lemma): L02 (LLC across severance) + L05-L06 (KMS state) $\Rightarrow$ the Jacobson thermodynamic derivation \cite{Jacobson1995}.  Status is AR because the $l_{\rm RA}$ identification (O06) is not yet a derived result; when O06 closes, D10 upgrades to DR automatically.

\section{Quantum Information: The Kinematic Coherence Bound}\label{sec:kcb}

\begin{leanbox}
\textbf{Theorem~L12 (Qubit Fragility) \statusLV:} Electrons and photons have minimum BDG stability score~1.  (\texttt{structural\_fragility}, RA\_D1\_Proofs.lean.)
\end{leanbox}

Each qubit is a kinematic site for spurious actualization.  The cumulative probability of spurious actualization for $N$ qubits:
\begin{equation}
  P_{\rm spurious}(N) = 1-(1-p_{\rm th})^N \approx N\cdot p_{\rm th}
\end{equation}
Setting $P_{\rm spurious}=\eta$ (quality factor) gives:
\begin{equation}
  N_{\max} = \eta\cdot p_{\rm th}
  \label{eq:KCB}
\end{equation}
Beyond $N_{\max}$, cumulative spurious actualization exceeds the fault-tolerance budget regardless of single-qubit fidelity.  Standard QM predicts exponentially difficult but unbounded scaling; RA predicts a hard ceiling.

This is the most near-term falsifiable prediction with strong discriminating power between RA and standard QM.

\section{Complete Predictions Table}\label{sec:predictions}

\begin{table}[ht]
\centering
\caption{All RA predictions, sorted by experimental proximity.  See Papers~I--IV for derivations.}
\label{tab:predictions}
\begin{tabular}{p{0.5cm}p{5.5cm}p{1.8cm}p{5cm}}
\toprule
ID & Prediction & Status & Experiment/Timeline \\
\midrule
P01 & BMV null; $\tau_{\rm act}\approx1.9\times10^4\,\mathrm{s}$ at $1\,\mathrm{mK}$ & \statusDR & Multiple groups pursuing; near-term \\
P02 & $H_0\propto\rho_b^{-1/2}(\text{LOS})$; $75.9\,\mathrm{km/s/Mpc}$ & \statusCV & DESI; current data \\
P03 & No DM signal below neutrino floor (categorical) & \statusDR & LZ/XENONnT; ongoing \\
P04 & $t^*=0.274/g$ (spin-bath, parameter-free) & \statusCV & SC circuits; medium-term \\
P05 & $N_{\max}=\eta p_{\rm th}$ (QEC hard ceiling) & \statusAR & QEC at $\sim10^3$ qubits; near-term \\
P06 & Life wherever F1+F2; substrate-independent & \statusAR & SETI biosignature search; long-term \\
P07 & $T^{-4}$ vs $T^{-1}$ decoherence scaling & \statusDR & Levitated nanoparticle experiments \\
P08 & No proton decay (baryon number topological) & \statusDR & Any future proton decay experiment \\
P09 & $\theta_{\rm QCD}=0$ exactly; no axion & \statusDR & Strong CP experiments; dark matter axion searches \\
\bottomrule
\end{tabular}
\end{table}

\section{Open Problems with Proof Paths}\label{sec:open}

\begin{table}[ht]
\centering
\caption{Open problems with current best paths to resolution.}
\label{tab:open}
\begin{tabular}{p{1cm}p{5cm}p{7cm}}
\toprule
ID & Problem & Path to closure \\
\midrule
O03 & Bianchi on curved background & Perturbative $(l_P/R)^2$ expansion: flat-space term (D04) + curvature correction bounded by LLC range (L01+L02). No Tomita-Takesaki needed at leading order. \\
O06 & $l_{\rm RA}$ from first principles & Dominant $k=2$ term in D4U02 = minimal BDG diamond; area $=l_P^2\Rightarrow l_{\rm RA}=l_P$; gives $\xi=\Delta S^*/l_P^2=C03/l_P^2$. \\
O10-G & Lean: $G_{\mu\nu}$ derivation & RA-native route: D04 + O10-spin2 directly; bypasses Weinberg S-matrix theorem. \\
O11 & Lorentz invariance from DAG & BDG modular invariants ($c_k\bmod\gcd$) as discrete Lorentz-invariants; question for Dowker collaboration. \\
O-I-1 & Wyler fractional correction & Derive group volume ratio from BDG dynamics. \\
O-I-2 & CKM matrix & Cabibbo angle conjecture $|V_{us}|=\sin(2/9)$; full quark mixing from SU(3)$_{\rm gen}$. \\
O-II-1 & Born rule derivation & Show $a(v|C)\propto\psi(v)$ follows from BDG action. \\
O-III-3 & Dark energy alternative & Quantify intrinsic graph growth rate as $\Lambda=0$ cosmological expansion mechanism. \\
\bottomrule
\end{tabular}
\end{table}

\section{Integrity Checks}\label{sec:ic}

The following integrity checks (IC) document known scope caveats and potential weaknesses.  They are recorded here not as defects to be hidden but as precise statements of what the framework currently claims and what it does not.

\textbf{IC01:} L04 uses $\sigma_0\propto I$ (finite-dimensional truncation).  The physical claim requires Poincar\'e invariance of the Minkowski vacuum, which is a standard QFT result but not the statement proven in Lean.

\textbf{IC02:} D02 (Unruh resolution) covers modular flow stationarity.  The explicit detector click story (no actualization in Minkowski vacuum) is now DR via O05r, which uses L05+timescale formula.

\textbf{IC03:} $l_{\rm RA}=\sqrt{4\Delta S^*}\ell_P$ is derived from area-law self-consistency (D08), not from BDG geometry independently (O06).

\textbf{IC04:} A01 (GR uniqueness) uses $d=4$ imported from L11 via RATM/RASM.  Should be derived directly in the Paper~III context.

\textbf{IC05:} A03 (BMV null mechanism) uses ``unactualized $\Rightarrow \rho_A=0$'' which requires O07 for the macroscopic COM case.  O07 is now DR, so this is substantially closed.

\textbf{IC06:} C01 ($\alpha_s$ result) gives the UV fixed point at $\mu\to\infty$.  The IR fixed point $\alpha_s=1/3$ is also proved \statusCV.  The relationship between the two is the BDG two-state RG flow.

\textbf{IC07:} L07 (causal invariance) was conditional on amplitude locality in RA\_AQFT\_Proofs\_v10.lean.  O01-O02 in RA\_AmpLocality.lean are unconditional.  L07 is superseded by O01+O02.

\section{The Philosophical Synthesis}\label{sec:philosophy}

\subsection{What Relational Actualism Claims About Physics}

Physics has operated for four centuries with a foundational metaphor: laws on one side, phenomena on the other.  The laws are fixed, background structure; the phenomena play out against them.  This structure is so deeply embedded that it is rarely stated explicitly.  It is simply how physics is done.

Relational Actualism dissolves this distinction.  The Local Ledger Condition (L01) is not a constraint that actualization events obey.  It \emph{is} what actualization events are, stated as a bookkeeping property.  Conservation of energy is not imposed on the causal graph from outside --- it is the LLC read at the appropriate scale.  General relativity is not a law governing spacetime --- it is the effective description of the actualization density field at scales far above the Planck length.

The arrow of time is not explained by the framework.  It is what the framework \emph{is}.  Every piece of RA mathematics is time-asymmetric because actualization is irreversible by definition.

\subsection{The Web, Not the Chain}

The dependency structure of RA is not a derivation chain (axioms to theorems in sequence).  It is a web.  The LLC (L01) connects to gravitational conservation (D04), quantum information (L03), biological complexity (via D11), and black hole thermodynamics (via D10).  The BDG threshold $\Delta S^*$ appears in quantum measurement, BMV experiments, biological chemistry, the Unruh effect, the Big Bang, and the origin of life --- not as an approximate connection but as the same formula evaluated at different scales.

The open problems (Table~\ref{tab:open}) are not independent mysteries.  Each is the same question --- ``when does a process become irreversible?'' --- asked at a different scale.  The answer is always $\Delta S^*=0.60069$, because that question has only one answer when irreversibility is the primitive.

\subsection{Coherence as Evidence}

The internal coherence of RA --- the fact that results at different scales connect through the same constants and the same causal graph structure --- is itself evidence of something.  It does not prove the framework is correct.  A beautifully coherent framework can be wrong, and RA makes specific predictions that can falsify it.

But coherence is evidence.  When a single commitment (irreversibility is primitive) simultaneously addresses the measurement problem, the cosmological constant, the Hubble tension, the origin of biological complexity, the number of spacetime dimensions, and the values of the coupling constants --- without free parameters and without adjusting the framework to fit --- this is not a coincidence.  It is either the beginning of a correct understanding of Nature, or it is the most instructive wrong answer in the history of physics.

Either outcome advances physics.

\section*{Acknowledgments}

This programme was developed over twenty-three years and formalised over approximately twelve months with AI collaboration (Claude, Anthropic; Gemini, Google).  The Lean~4 infrastructure builds on the Lean-QuantumInfo library (Meiburg, Perimeter Institute).  All papers are open-access at \url{https://zenodo.org/communities/relational-actualism}.  Lean~4 files, Python scripts, and figure sources are at \url{https://github.com/jsandeman/Relational-Actualism}.

\bibliographystyle{plain}
\begin{thebibliography}{99}
\bibitem{RA_P1} Sandeman, J.~F. (2026). Paper~I. Zenodo, doi:10.5281/zenodo.19362289.
\bibitem{RA_P2} Sandeman, J.~F. (2026). Paper~II. Zenodo, doi:10.5281/zenodo.19362968.
\bibitem{RA_P3} Sandeman, J.~F. (2026). Paper~III. Zenodo, doi:10.5281/zenodo.19363017.
\bibitem{RA_P4} Sandeman, J.~F. (2026). Paper~IV. Zenodo, doi:10.5281/zenodo.19198224.
\bibitem{Dowker2013} Dowker, F., and Glaser, L. (2013). \textit{CQG}, 30, 195016.
\bibitem{Jacobson1995} Jacobson, T. (1995). \textit{Phys.\ Rev.\ Lett.}, 75, 1260.
\bibitem{Weinberg1964} Weinberg, S. (1964). \textit{Phys.\ Rev.}, 135, B1049.
\bibitem{Bombelli2009} Bombelli, L., Henson, J., and Sorkin, R.~D. (2009). \textit{Mod.\ Phys.\ Lett.~A}, 24, 2579.
\bibitem{Smolin1992} Smolin, L. (1992). \textit{CQG}, 9, 173.
\bibitem{Mathlib} Mathlib Contributors (2024). Lean~4. \url{https://leanprover-community.github.io/}
\end{thebibliography}

%═══════════════════════════════════════════════════════════════════════════════
\section{Integrity Checks: Known Scope Caveats}\label{sec:ic}
%═══════════════════════════════════════════════════════════════════════════════

The following integrity checks (IC) document precisely what the framework currently claims and what it does not.  They are recorded not as defects but as honest statements of scope.  Every theoretical programme has limits; the RA programme documents them explicitly.

\subsection{Layer~0 Integrity Checks (Lean-Verified)}

\paragraph{IC01 (L04 scope).}
Theorem~L04 (Frame Independence) is proved in a finite-dimensional truncation with $\sigma_0\propto I$.  The physical claim requires Poincar\'e invariance of the Minkowski vacuum, which is a standard result in algebraic QFT (Haag \cite{Haag1996}).  The gap: finite-dimensional $\sigma_0\propto I$ is not Poincar\'e-invariant (it is invariant under all unitaries, not just Poincar\'e ones).  The physical significance is limited: the finite-dimensional result is sufficient for the main application (Lorentz covariance of the actualization criterion, D01), and the extension to the full AQFT setting is expected to hold.

\paragraph{IC07 (L07 status).}
L07 (conditional causal invariance) is superseded by O01+O02 (unconditional, any permutation).  The RA AQFT files contain both.  For any paper that cites causal invariance, the unconditional version (O01+O02 from RA\_AmpLocality.lean) should be used.

\subsection{Layer~1 Integrity Checks (Computation-Verified)}

\paragraph{IC06 ($\alpha_s$ fixed points).}
The computation C01 establishes the UV fixed point $\alpha_s(m_Z)=1/\sqrt{72}$ at $\mu\to\infty$.  The IR fixed point $\alpha_s=1/3$ at $\mu\to0$ is also computation-verified but uses the two-state model approximation.  Papers discussing the strong coupling should note both fixed points.

\paragraph{IC-$\Delta S^*$ (Irrationality).}
$\Delta S^*=0.60069$ is established as the exact value to $10^{-10}$ precision.  The rational approximation $3/5=0.600$ is ruled out at $25.2\sigma$.  Any claim that uses $\Delta S^*$ numerically should use $0.60069$ (or the more precise value $0.548436...$ for $P_{\rm acc}(1)$), not $3/5$.

\subsection{Layer~2 Integrity Checks (Derived)}

\paragraph{IC02 (D02 gap).}
D02 (Unruh resolution via modular flow stationarity) covers the statement that $\Delta S_{\rm per\,photon}=0$ for the Rindler bath.  The full story of ``no detector clicks'' (O05r) requires the additional step of applying the universal timescale formula.  Papers presenting D02 should cite O05r for the complete argument.

\paragraph{IC03 ($\xi$ provisional).}
$\xi=1/(4\ell_P^2)$ is derived from the area-law self-consistency $\ell_{\rm RA}=\sqrt{4\Delta S^*}\ell_P$.  This is a self-consistency condition, not an independent derivation.  Until O06 is resolved, papers should present $\xi$ as ``provisionally derived'' rather than ``derived from first principles.''

\paragraph{IC04 ($d=4$ importation).}
A01 (GR uniqueness via Lovelock) uses $d=4$ imported from L11 via RATM/RASM.  The direct derivation of $d=4$ in the GR context (without reference to the SM topology papers) has not been written.

\paragraph{IC05 (Macroscopic COM actualization).}
A03 (BMV null mechanism: superposed mass sources zero metric) uses ``unactualized $\Rightarrow \rho_A=0$''.  For a single quantum particle, this is straightforward.  For a macroscopic center-of-mass in spatial superposition, the argument requires that $\tau_{\rm act}\gg\tau_{\rm exp}$ for the COM --- which is established by O07r.  Papers presenting A03 should cite O07r as the timescale closure.

%═══════════════════════════════════════════════════════════════════════════════
\section{The Philosophical Synthesis: Laws Are Phenomena}\label{sec:philosophy_full}
%═══════════════════════════════════════════════════════════════════════════════

\subsection{The Standard Picture and Its Hidden Assumption}

Physics operates with an implicit hierarchy: laws govern phenomena.  The laws of nature are the fixed, background structure.  The phenomena --- particle collisions, chemical reactions, biological processes --- play out against this background.  This hierarchy is so deeply embedded that it is rarely stated: it is simply how physics is done.

The laws are taken to have a different ontological status from the phenomena.  They are not observed; they are inferred.  They are not events; they are regularities.  A law says something like ``whenever event A occurs, event B will follow'' --- but the law itself is not an event.  It is a relation between events, a description of the structure of reality, not a part of reality itself.

This hierarchy has a name in philosophy of science: it is the nomological account of laws (from the Greek \textit{nomos}, law).  Nature has laws; the laws govern nature; the laws are not themselves part of nature.

\subsection{How RA Dissolves the Hierarchy}

Relational Actualism makes the nomological account mathematically untenable.

The Local Ledger Condition (Theorem~L01) is proved in Lean~4.  It states: at every vertex of the causal DAG, $\Sigma_{\rm out}(v)=\Sigma_{\rm in}(v)$.  Is this a law or a phenomenon?

It is not a law in the nomological sense: it does not say ``whenever an event occurs, this balance holds.''  It is a property of every event: it is what every actualization event \emph{is}, stated as a bookkeeping property.  An actualization event with $\Sigma_{\rm out}\neq\Sigma_{\rm in}$ is not a violation of a law --- it is a self-contradiction.  It would not be an event in the RA sense at all.

The LLC is not a constraint imposed on events from outside.  It is part of the definition of what an event is.

The same is true throughout RA:
\begin{itemize}
  \item Conservation of energy is the LLC at the appropriate scale --- not a law imposed on the causal graph but a property of every vertex in it.
  \item The arrow of time is not explained by RA --- it is what RA \emph{is}.  Actualization is irreversible by definition.  There is no RA formalism that runs backwards.
  \item The metric is not evolved by the Einstein equations \emph{as a law}; it is what the actualization density field \emph{looks like} at macroscopic scales.  The Einstein equations are the effective description of this appearance.
\end{itemize}

\subsection{The Web Is the Theory}

The dependency graph of RA is not a derivation chain (axioms at top, theorems at bottom, flowing in one direction).  It is a web.  Every node has multiple parents and multiple children.  The LLC (L01) connects to gravitational conservation (D04), quantum information (L03), biological complexity (D11), and black hole thermodynamics (D10).  The threshold $\Delta S^*$ appears in quantum measurement, BMV experiments, biological chemistry, the Unruh effect, the Big Bang, and the origin of life.

Any attempt to read the web as a chain --- ``first we prove the LLC, then the BDG action, then the field equations, then the predictions'' --- misses the structure.  The chains exist (the dependency graph is acyclic, so there are chains), but they are not the theory.  The theory is the web: the set of all the connections, the way the same fact appears in fourteen different physical contexts, the way $\Delta S^*$ is simultaneously the measurement threshold and the origin-of-life threshold and the universe-seed threshold.

The mutual implication structure \cite{RA_MutualImplication} is the correct framing.  The laws and the phenomena are not in a hierarchy.  They are in mutual implication: the laws are the phenomena viewed at the level of pattern, and the phenomena are the laws instantiated at specific events.  Neither is more fundamental than the other.

\subsection{What It Would Mean for RA to Be Wrong}

A beautiful framework can be wrong.  RA makes specific predictions that can falsify it:
\begin{itemize}
  \item If the BMV experiment detects gravity-mediated entanglement, the core claim that superposed mass sources zero metric is false.
  \item If DESI finds no correlation between $H_0$ and baryon density along the line of sight, the Hubble tension mechanism is wrong.
  \item If LZ or XENONnT detects a WIMP, the categorical exclusion of WIMPs fails.
  \item If a fourth generation of quarks or leptons is discovered, the BDG closure theorem is wrong.
  \item If proton decay is observed, exact baryon conservation fails.
\end{itemize}

Each of these is independently falsifiable.  The framework does not protect itself from falsification by post-hoc adjustment.  If any prediction fails, the failure is a precise statement of what aspect of the framework is wrong, and it points directly to the physics that needs to change.

\subsection{The Significance of Internal Coherence}

The absence of free parameters is not merely an aesthetic virtue.  It makes the framework's coherence epistemically significant.

If a framework has $N$ free parameters, any agreement between predictions and observations could be attributed to parameter tuning.  The framework would be consistent with a wide range of observations, and its coherence would be purchased by fitting.

RA has zero free parameters for any of the predictions in Table~\ref{tab:predictions}.  Every number --- $\alpha_{\rm EM}^{-1}=137$, $\alpha_s=0.11785$, $H_0^{\rm Eridanus}=75.9$, $f_0=5.42$, $t^*=0.274/g$ --- is a consequence of the same five integers that were fixed by the geometry of the d'Alembertian.  The coherence is not purchased by fitting.  It is either genuine --- the framework correctly captures something about the structure of nature --- or it is an extraordinary accident.

The probability that fourteen independent predictions all agree with experiment by accident, in a framework with zero free parameters, is essentially zero.  The coherence is evidence.

\end{document}
